\documentclass[11pt,a4paper]{article}
\usepackage[utf8]{inputenc}
\usepackage[T1]{fontenc}
\usepackage{geometry}
\geometry{margin=1in}
\usepackage{hyperref}
\usepackage{enumitem}
\usepackage{booktabs}
\usepackage{longtable}
\usepackage{xcolor}
\usepackage{titlesec}
\usepackage{amsmath,amssymb}

\titleformat{\section}{\Large\bfseries\color{blue!70!black}}{}{0em}{}
\titleformat{\subsection}{\large\bfseries\color{blue!50!black}}{}{0em}{}
\titleformat{\subsubsection}{\normalsize\bfseries\color{blue!30!black}}{}{0em}{}

\title{\textbf{Replication Plan}\\[0.5em]
\large A Portfolio Approach to Massively Parallel Bayesian Optimization\\[0.3em]
\normalsize OSTI ID: 2571540}
\author{Auto-generated Replication Blueprint}
\date{\today}

\begin{document}
\maketitle
\tableofcontents
\newpage

%---------------------------------------------------------------------
\section{Paper Overview}
%---------------------------------------------------------------------
This paper proposes a portfolio-based strategy for massively parallel Bayesian optimization (BO). Instead of relying on a single acquisition function, the method maintains a portfolio of diverse acquisition functions (e.g., Expected Improvement, Upper Confidence Bound, Thompson Sampling, Knowledge Gradient) and allocates the parallel batch budget across them using a multi-armed bandit or portfolio optimization framework (e.g., based on the Hypervolume Sharpe Ratio). This approach provides robustness against the ``exploration--exploitation'' dilemma and enables efficient use of hundreds to thousands of parallel evaluations.

%---------------------------------------------------------------------
\section{Detailed Workflow}
%---------------------------------------------------------------------

\subsection{Phase 1: Bayesian Optimization Framework}
\begin{enumerate}[label=\textbf{1.\arabic*}]
  \item \textbf{Implement or use a GP-based Bayesian optimization library.}
    \begin{itemize}
      \item Gaussian process surrogate model with specified kernel (e.g., Mat\'{e}rn 5/2).
      \item Hyperparameter optimization via marginal likelihood.
    \end{itemize}
  \item \textbf{Implement individual acquisition functions.}
    \begin{itemize}
      \item Expected Improvement (EI).
      \item Upper Confidence Bound (UCB) with varying $\beta$.
      \item Thompson Sampling (TS).
      \item Knowledge Gradient (KG).
      \item Probability of Improvement (PI).
    \end{itemize}
\end{enumerate}

\subsection{Phase 2: Portfolio Allocation Strategy}
\begin{enumerate}[label=\textbf{2.\arabic*}]
  \item \textbf{Implement the portfolio allocation mechanism.}
    \begin{itemize}
      \item Multi-armed bandit approach (e.g., Exp3, Hedge).
      \item Hypervolume Sharpe Ratio (HVSR) for multi-objective allocation.
      \item Uniform allocation baseline.
    \end{itemize}
  \item \textbf{Batch generation.}
    \begin{itemize}
      \item For a batch of $B$ parallel evaluations, allocate $B_i$ to acquisition function $i$.
      \item Generate candidate points by optimizing each acquisition function.
      \item Deduplicate and ensure diversity in the batch.
    \end{itemize}
\end{enumerate}

\subsection{Phase 3: Benchmark Functions}
\begin{enumerate}[label=\textbf{3.\arabic*}]
  \item \textbf{Implement standard benchmark functions.}
    \begin{itemize}
      \item Branin, Hartmann-6D, Ackley, Rosenbrock, Levy, Styblinski-Tang.
      \item Higher-dimensional test functions as specified.
    \end{itemize}
  \item \textbf{Implement real-world surrogate benchmarks} (if used in the paper).
\end{enumerate}

\subsection{Phase 4: Experimental Runs}
\begin{enumerate}[label=\textbf{4.\arabic*}]
  \item \textbf{Run the portfolio BO} on each benchmark for multiple random seeds.
  \item \textbf{Run baselines}: single-acquisition-function BO (EI-only, UCB-only, TS-only), random search.
  \item \textbf{Vary batch sizes} (e.g., $B = 8, 16, 32, 64, 128, 256$).
  \item \textbf{Record}: best observed value vs.\ iteration, wall-clock time, portfolio weights over time.
\end{enumerate}

\subsection{Phase 5: Analysis and Validation}
\begin{enumerate}[label=\textbf{5.\arabic*}]
  \item \textbf{Compute performance metrics}: simple regret, cumulative regret, normalized improvement.
  \item \textbf{Statistical comparison}: mean and confidence intervals over random seeds.
  \item \textbf{Reproduce all figures and tables} from the paper.
  \item \textbf{Analyze portfolio weight dynamics.}
\end{enumerate}

%---------------------------------------------------------------------
\section{Datasets}
%---------------------------------------------------------------------

\begin{longtable}{p{4.5cm} p{3.5cm} p{6cm}}
\toprule
\textbf{Dataset / Source} & \textbf{Type} & \textbf{Location / Notes} \\
\midrule
\endfirsthead
\toprule
\textbf{Dataset / Source} & \textbf{Type} & \textbf{Location / Notes} \\
\midrule
\endhead
Benchmark test functions & Synthetic & Standard definitions; implemented programmatically \\
\midrule
Real-world surrogates (if used) & Public & HPOBench, Bayesmark, or as specified \newline \url{https://github.com/automl/HPOBench} \\
\bottomrule
\end{longtable}

%---------------------------------------------------------------------
\section{Models, Tools, and Software}
%---------------------------------------------------------------------

\begin{longtable}{p{4.5cm} p{2.5cm} p{6.5cm}}
\toprule
\textbf{Tool / Library} & \textbf{Version} & \textbf{Purpose / URL} \\
\midrule
\endfirsthead
\toprule
\textbf{Tool / Library} & \textbf{Version} & \textbf{Purpose / URL} \\
\midrule
\endhead
BoTorch & $\geq 0.8$ & Bayesian optimization in PyTorch \newline \url{https://botorch.org/} \\
\midrule
GPyTorch & $\geq 1.9$ & Gaussian process modeling \newline \url{https://gpytorch.ai/} \\
\midrule
PyTorch & $\geq 1.12$ & Deep learning backend \\
\midrule
Ax Platform & latest & High-level BO framework (uses BoTorch) \newline \url{https://ax.dev/} \\
\midrule
NumPy / SciPy & $\geq 1.21$ & Numerical operations, optimization \\
\midrule
Matplotlib / Seaborn & $\geq 3.4$ & Performance comparison plots \\
\midrule
scikit-optimize & latest & Alternative BO library \newline \url{https://scikit-optimize.github.io/} \\
\bottomrule
\end{longtable}

\subsection{Hardware Requirements}
\begin{itemize}
  \item Multi-core workstation for parallel batch evaluation.
  \item GPU optional (useful for GP fitting with GPyTorch on large datasets).
  \item Each benchmark run: minutes to hours depending on batch size and function evaluation cost.
\end{itemize}

\end{document}
